\section{Patient-Level Clustering Extension}\label{sec:extension_methods}

An exploratory follow-up to this study was conducted on October 15, 2026. The analysis explored whether the existing historical severity categories represent stable subcategories with regard to different ways the data can be represented, with regard to how the data may have been missed, and when using different sampling units. The findings from this study are supplemental to those found in H1-H4, but do not redefine those hypotheses. The supporting literature was limited to journal articles published between October 15, 2025 and October 15, 2026. Articles used in comparisons of mixture modeling of missing data, multi imputation ensemble approaches and techniques for diagnosing clusterability (in addition to comparative longitudinal analysis) included \textcite{tong2025mixturemissing,tomo2026ensemble,neville2026clusterability,julia2026longitudinal} and provided further rationale for determining if there existed structure in the measure feature space based upon distortions present in UMAP embeddings \parencite{chen2026umapartifacts}. None of these articles validated a particular type of ABI or established which methodology would produce better results than the historical methodology.

\subsection{Cohort and Discovery--Validation Split}

We identified an eligible population of 17,406 evaluations performed on 7,285 individuals, across the 14 retained cognitive tests. Identifiers for patients link repeated evaluations; there are no duplicated patient-evaluation-time combinations. No validated identifiers exist for defining episodes of injury. One patient reported differing ages at injury across assessments and was removed from trajectory modeling based on their evaluation history.

Patient assessments were ordered by time post-injury, and each patient's first eligible evaluation was used in the cross-sectional aspect of this study. When determining the order of assessments for multiple evaluations within a single patient, source row ordering is used to break ties deterministically. In a fixed random split using seed \#20261015, 5,828 patients were randomly assigned to the training set (discovery) and 1,457 patients were randomly assigned to the testing set (validation). Since each patient provides only one index profile, no patient appears in both groups. This allows us to estimate structure in patient index profiles while we describe assessment state across follow-up in the historical analysis. It has also already been noted that the 14 feature variables chosen for use in this research were those same as the feature variables used in our previous historical analyses for the entire cohort of patients. Thus, the division as described does not represent an independent validation of how we chose these feature variables.

Eligibility required at least one observed retained test. A separate sensitivity analysis evaluated the impact of sparse measurements of patient assessment profiles. Specifically, it used index evaluations where at least seven of the 14 retained tests had been measured.

\subsection{Shared Preprocessing and Feature Representations}

Data in discovery datasets establish the 1st and 99th percentiles, means and standard deviations for each test. Values are clipped at those percentile bounds, and timed \texttt{ATMTA} scores are sign-inverted. Standardization is employed to create larger values associated with higher achievement. This transformation leaves missing entries unchanged. For completed-data branches, a bounded Bayesian chained-equations imputer was fitted on discovery patients and used to impute missing values for validation patients. An iterative, deterministic reference model was created using 15 iterations. An additional sensitivity study utilizing 50 iterations was conducted to assess convergence of the iterative algorithm. Note that this imputation method differs from the historical \texttt{miceforest} implementation and has its own modeling assumptions.

The completed standardized tests were grouped by six domains based on the frequency of each test's administration. Each domain's weight is determined by the frequency of each test being administered to the discovery dataset; however, the weights have been normalized within each domain. Domain median and inter-quartile range (robust scaling) derived from the discovery dataset were utilized to calculate the reference distance representation. The comparator options include either the 14 completed tests or the 5 non-orientation domains. Following removal of the first principal component derived from the discovery patient data, a five-dimensional PCA basis was developed. Residual representation represents a diagnostic tool for identifying potential shape differences amongst profile groups. Therefore, residual representation may remove valid signals as well as common variability.

After examining discovery-set variances, an additional scaling sensitivity was introduced in which each domain has a variance of one. This sensitivity is designed to evaluate the influence of domains having "ceiling" issues which tend to have very narrow ranges around their midpoints. By evaluating identical label sets in both geometric models, we can separate differences in how the data are divided from differences in how silhouette is calculated. The transformations applied to the validation patients remain unchanged from those originally estimated in the discovery data.

In terms of analyzing each patient individually, the patient level analysis utilizes clipping, direction alignment and test standardization before the values are then combined. Previous supplementary analyses utilized multiple imputation and complete case comparison methodologies to aggregate raw tests. Consequently, the prior results that were provided as part of this study used a different preprocessing methodology than the patient-level extension. Controlled comparisons using identical completed test values assess the impact of aggregation without recalculating all the previous H1-H4 statistics.

\subsection{Original-Space and Incomplete-Data Models}

The clustering techniques k-means and PAM (sampled partitioning around medoids) were evaluated for values of $k=2,\ldots,8$. PAM utilizes CLARA-style sampling on three sets of 400 discovery patient data, followed by a BUILD/SWAP optimization within each sample. Additionally, the Gaussian Mixture Model was evaluated for $k=1,\ldots,8$. Both full and tied covariance matrices were used. HDBSCAN was run against the data for two sets of parameters: (a) minimum cluster size of either 100 or 300; and (b) \texttt{min\_samples}=20. For the comparison utilizing UMAP, three-dimensional representations were employed, along with 15 or 50 nearest neighbor values with a \texttt{min\_dist}=0 before applying the same HDBSCAN parameters as above. In terms of assigning validation patients to clusters, this was performed using the fitted centers, medoids, mixture probabilities, or UMAP transformation and HDBSCAN approximate prediction.

Models to fit each representation/method were selected based upon the highest discovery silhouette for partition/density methods and lowest discovery BIC for Gaussian mixtures. When applicable, fits had to converge and every group had to contain at least two percent of the discovery patients. Silhouettes were computed in the respective original representation (including UMAP-derived labels). Up to 2,000 randomly sampled observation discovery scores were used when computing discovery scores. All available assignments were used when calculating validation scores. Noise was excluded from silhouette, and coverage was reported separately. Scores from different feature geometries are not a common performance scale. A silhouette search for $k\geq2$ does not determine whether there exists a group relative to an existing one population model.

The incomplete-data mixture models were fitted to 14 transformed test results using expectation maximization (EM). Missing coordinates were integrated out in the observed-data likelihood. There were three restricted families used: (a) Gaussian; (b) $t$-distributions with df = 4 \& 10; (c) contaminated normal with a 90\% normal core and tenfold covariance inflation for the remaining 10\%. Each of the above mentioned families were estimated with $k=1,\ldots,8$, using two different starting conditions, with a 120-iteration limit or convergence when the change in log likelihood per observation was below $10^{-5}$, and an eigenvalue floor of 0.02 to the dispersion matrix. Due to scale matrices being used within the $t$ family estimation process, the eigenvalue floor also served as a scale-matrix floor. Mean Imputation was the only method used for developing starting values. Verification of these implementations has been done through computational verification. These are simplified versions of the full \textit{MixtureMissing} package \parencite{tong2025mixturemissing} and lack some functionality.

BIC, entropy-penalized ICL, convergence, validation observed-data log likelihood and posterior membership probabilities were recorded for each fit. Selected fits had to converge and assign at least 2\% discovery membership to each component. Likelihood measures are compared only among fits sharing the same transformation into the 14-dimensional space. Because continuous distributions are being used to approximate discrete/bounded scores, any additional groups could represent ceiling effects, skewness, or tailing rather than separate subpopulations of patients.

\subsection{Independent Refitting, Consensus and Cluster Tendency}

Twenty stochastic imputed datasets were generated from bounded Bayesian posterior draws, using ten chained-equations iterations per dataset. K-means and sampled PAM were fitted separately with fixed $k=2,3,4$, refitting domain scaling for each draw. Ensemble PAM used two samples of 300 discovery patients. Each assignment vector contained the same 1,457 validation patients. Agreement across draws was measured by mean pairwise ARI and empirical 5th-95th percentiles. These are not confidence intervals because partition pairs share fits. A separate discovery-silhouette k-means screen assessed $k=2,\ldots,8$ in each run.

Fifty runs used 80\% subsamples of discovery patients to fit preprocessing, imputation, scaling, and the fixed k=3 group k-means. Fifty additional runs were made where the clustering seed varied while the input vector was kept constant. The patient sampling variability was assessed by computing the held-out ARI and matched per-cluster Jaccard indices to differentiate it from initial condition variability. Additional sensitivity checks assessed the effects of mean versus KNN imputation, no clipping, stricter availability requirements for tests, and independent fitting of the validation partitions.

Finally, I compared average-linkage co-assignment, greedy mean-NMI, and generic NMF as three approaches to summarize the descriptive properties of held-out assignment vectors using methods similar to \textcite{tomo2026ensemble}. Although these are the native implementations of those methods, they are not an exact replication of that paper. Additionally, consensus regarding the validation assignments does not constitute a predictive model which has been separately validated. Collapsed groups and nonconvergence are reported.

The results from the PCA/DIP diagnostic analysis were then compared to a matched-size screening set of 1,000 discovery patients and 100 separate reference samples that were each constructed by sampling from a single Gaussian copula. The empirical discrete marginal distributions as well as the missingness joint mask created through the application of resampling techniques were preserved for each. Additionally, each of these 100 reference datasets were subjected to the exact same processing procedures used on the initial dataset. Specifically, this included: preprocessing; estimation of missing data; and K-Means Clustering Analysis with varying cluster counts ($k=2,\ldots,8$). The tail fraction comparisons include both the observed Max Silhouette and PC1 Dip Statistics when compared to those present within the reference distribution(s). The method of comparison described above is known as ``descriptive fitted-reference calibration''. The simulation missingness masks are generated independently of the simulated data and selection across all model families is not repeated. In addition, there is a two group positive control check to check basic recovery. However, neither of the screens or the projection based modality test identify clinical subtypes \parencite{neville2026clusterability}.

\subsection{Patient Trajectories and Reproducibility}

To qualify for the longitudinal sample, a patient had to meet all of the following requirements. There had to be at least three visits with seven or more observed tests and at least three distinct time points. The interval between the first and last eligible observation could be no shorter than 0.25 years and no longer than five years. Recorded age at injury also had to be consistent across assessments.

I developed a longitudinal cognitive score using the frozen index-discovery preprocessing, imputer and robust-space PC1. The score was oriented towards better performance and divided by its index-discovery standard deviation. Time was measured in years since the first eligible time point. Each patient kept their original discovery or validation status throughout the longitudinal analysis.

I modelled the data using Gaussian linear trajectory mixture models. The patients' random intercepts and slopes were integrated into the models. Models varying in number of classes (from 1 to 4) were fit. Random-effects covariance and residual variance were shared across all of the classes. For each model, three different sets of initial values were tested and up to 300 EM iterations were allowed. Only models that converged as well as had an average posterior probability of at least 0.05 in each of their respective classes could be evaluated with respect to discovery BIC. Therefore, any model that failed to converge during its estimation procedure or had an average posterior probability < 0.05 in any particular class would be eliminated from consideration. To provide a point of comparison, we also ran a one-class model in which there was a continuous random effect. Additionally, patient subsampling was conditional on selecting more than one class. Rather than determining how much of the identified pattern(s) could be due to underlying causal mechanism(s), we wanted to determine what other observed patterns of interest existed. We did not run analyses of nonlinear trajectories. Also, we did not utilize the recently developed multivariate contaminated longitudinal model. Recent studies of longitudinal data motivate separating initial value from rate of change \parencite{julia2026longitudinal}.

All analyses were performed using Python 3.12.14 with fixed version numbers for all dependencies, reproducible random seed states and as previously stated, model grids. Eight numerical tests evaluated the reference likelihoods, recovery of missing data, monotonicity of EM log-likelihood, preprocessing directionality, selection of the medoids, integration over trajectory likelihoods and evaluation of the overall ensemble objective function. No independent rehabilitation outcome was available, so predictive clinical utility remains untested.
